\newpage 
\section{Relational Algebra}

\subsection{Relations}

A \textbf{relation} is a set of the Cartesian Product of two sets. It contains ordered pairs \((a,b)\).

The Cartesian Product for our entities is defined as

\[
    S \times T := \bigcup_{(s_1, \dots, s_n) \in S} \left[ \bigcup_{(t_1, \dots, t_k) \in T} \{(s_1, \dots, s_n, t_1, \dots, t_k)\}\right]
\]

Simplified each element of \(S\) is combined with each element of \(T\). The size of this
special product is \( |T \times S| = |S| \cdot |T| \).

In the relational algebra relations represent tables and their rows the tuples.

\subsection{Relational Algebra}

Relational Algebra is a \textbf{formal, procedural query language} for the relational model. It provides a set of \textbf{operators} that take relations (tables) as input and produce relations as output. It forms the \textbf{theoretical foundation of SQL} and relational database systems.

Throughout this section, we use two example relations to illustrate each concept.

\subsubsection{Axioms of Relational Algebra}

Relational Algebra is based on \textbf{set theory} and satisfies the following axioms:

\begin{itemize}
\item \textbf{Closure}: Every operation produces a relation.
\item \textbf{Set Semantics}: Relations do not contain duplicate tuples.
\item \textbf{Order Irrelevance}: The order of tuples and attributes does not matter.
\item \textbf{Composability}: Operations can be nested arbitrarily.
\end{itemize}

\textbf{Example:}
  
If \(\sigma_{id > 1}(A)\) returns a relation, it can be used as input to another operator such as projection.

\subsubsection{Relations}

A \textbf{relation} is a finite set of tuples over a fixed set of attributes.

Formally:

\[
    R \subseteq D_1 \times D_2 \times \dots \times D_n
\]

\textbf{Example:}
  
Relation \(A(id, value)\) is a set of tuples:

\[
    A = \{(1, X), (2, Y)\}
\]

\subsubsection{Identifier and Schema}

We define  

\[
    \mathrm{ident}(R) = \{a_i\}_i \quad \text{ and } \mathrm{schema}(R) = \{[a_1:T_1, \dots, a_n:T_n]\}
\]

\(\mathrm{ident}\) returns the names of the attributes of the relation, while 
\(\mathrm{schema}\) the attributes and its corresponding types.

\subsection{Domain of an Attribute}

The \textbf{domain} of an attribute \(a\) in a relation \(R\) is the set of all possible values that \(a\) can take in \(R\).

\subsubsection{Cartesian Product for Relations}

The \textbf{Cartesian product} combines every tuple of one relation with every tuple of another.

\[
S \times T = \bigcup_{(s_1, \dots, s_n) \in S} \left[ \bigcup_{(t_1, \dots, t_n) \in T} \{(s_1, \dots, s_n, t_1, \dots, t_k)\} \right]
\]

With size \(|S \times T|  = |S| \times |T|\).

\subsubsection{Union \(\cup\)}

Combines tuples from two relations with the \textbf{same schema}.

\[
A \cup B
\]

\subsubsection{Set Difference \(\setminus\)}

Returns tuples in one relation but \textbf{not} in another.

\[
A \setminus B
\]

\subsubsection{Intersection \(\cap\)}

Returns tuples common to both relations.

\[
A \cap B = \{\}
\]

\subsubsection{Selection \(\sigma\)}

Filters tuples based on a predicate.

\[
\sigma_{\theta}(S) = \{s \in S \mid s \text{ satisfies } \Theta\}
\]

\subsubsection{Projection \(\pi\)}

Selects specific attributes.

\[
\prod_{a_{i1}, \dots, a_{ik}}(S) \{(a_{i1}, \dots, a_{ik}) \mid a \in S\}
\]

\subsubsection{Symmetric Difference \(\triangle\)}

Tuples in either relation but \textbf{not in both}.

\[
A \triangle B = (A \setminus B) \cup (B \setminus A)
\]

\subsubsection{Rename \(\rho\)}

Renames relations or attributes.

\[
\rho_{A'(key, data)}(A)
\]

\subsubsection{\texorpdfstring{\(\mathrm{with}\)}{with}}

\(\mathrm{with}\) allows us to store products of specific operation with an 
special name which can be used later in oter Operations.

\subsubsection{Join \(\bowtie\)}

A \textbf{theta join} is the Cartesian Prouct plus a selection.

\[
    S \bowtie_{\theta } T = \sigma_\theta (S \times T)
\]

\paragraph{Inner Join \(\bowtie\)}

A \textbf{theta join} combines tuples that satisfy a condition.

\[
    S \bowtie_{\theta } T = \sigma_\theta (S \times T)
\]

\paragraph{Equi Join}

An equi-join is a join where the join condition consists only of equality 
comparisons between attributes of the two relations.

\[
    R \bowtie_{R.a_1 = S.b_1 \land \dots \land R.a_k = S.b_k} S
\]

The name and domain of the attribures must be the same.

\paragraph{Natural Join}

It is a speical case of the equi join which can only be performed only if there 
is a common attribute.

\begin{itemize}
\item Natural Join = Cartesian Product + Selection + Projection
\end{itemize}

\paragraph{Outer Join}

Preserves non-matching tuples.

\begin{itemize}
\item Left outer join 
\item Right outer join
\item Full outer join
\end{itemize}

\textbf{Example (Left Outer Join)}

\begin{tabular}{|c|c|c|}
\hline
id & A.value & B.value \\
\hline
1 & X & NULL \\
2 & Y & Y \\
\hline
\end{tabular}

\paragraph{Semi Join}

Returns tuples from one relation that \textbf{have a match} in the other.

\[
A \ltimes B
\]

\textbf{Example:}


\begin{tabular}{|c|c|}
\hline
id & value \\
\hline
2 & Y \\
\hline
\end{tabular}

\subsubsection{Division \(\div\)}

Used to express \textbf{"for all"} queries. It can only be performed when 
the relation on the right is a proper subset of the left relation.

Let:
\begin{itemize}
\item \(R(X, Y)\)
\item \(S(Y)\)
\end{itemize}

Then:

\[
R \div S
\]

\textbf{Example Concept}  
If \(A(id, value)\) contains all values in \(B(value)\), division returns the corresponding \texttt{id}s.

Division is commonly used in queries like:
\begin{quote}
“Find entities related to \textbf{all} items in a set.”
\end{quote}

\subsubsection{Assigment \((=)\)}

This is just using to store the result of a relational algebra operation.

\subsection{Types of Keys}

A \textbf{key} is an attribute or combination of attributes (\textbf{composite key}) which uniquely identifies an entity.

\begin{itemize}
\item \textbf{Candidate Key:} The term refers to the attributes which are suited for being the key.
\item \textbf{Irreducible:} A key has to be as atomic as possible.
\item \textbf{Super key:} It consists of a candidates key's attributes + potentially some extra attributes.
\item \textbf{Alternate Key:} This is an alternative primary key. This means that the UNIQUE constraints will be imposed on the attribute/s.
\item \textbf{Surrogate Key:} This key has only meaning inside the database. Example: ID \(\to\) Person's Name.
\item \textbf{Foreign Key} This key is an attribute of one table which references another table.
\end{itemize}

\subsection{Relation Schema}

The \((\mathcal{R}, \tau)\) represents a \textbf{relation schema}. We can constrain our relations to obey certain rules. Thus, limiting our space. Here, \(\mathcal{R}\) is the set of all relations of \(n\) elements and \(\tau\) is a \(n\) tuple (attributes).

This is the analog of using \verb|CREATE TABLE| in SQL. We are creating rules for 
the tuples inside the relation.

\subsection{Functional Dependencies}

A \textbf{functional relation} is dependence between one or more attributes commonly noted \texttt{A -> B}.

Formally \(a \rightarrow b\) holds:

\[
    \forall r (\mathcal{R}) \forall d_1, d_2 \in r (d_1[\alpha] = d_2[\alpha] \implies d_1[\beta] = d_2[\beta])
\]

where \(r(\mathcal{R})\) is a relation in the Schema \(\mathcal{R}\), \(\alpha\)
is a tuple in the schema and \(d[\alpha]\) is the tuple without \(\alpha\).

In simple terms the notation \( X \to Y\) means "we know Y if we know X".

A complete dependece is when \(A \to B\) is when we can not further reduce the tuple \(A\) 
so that we still get \(B\).

\subsection{Armstrong Axioms}

For a relation \(R = \{a_1, \dots, a_n\}\) with the schema \(\alpha, \beta, \gamma \subseteq \mathrm{ident}(R)\)
The Armstrong axioms are given: 

\begin{itemize}
\item \textbf{Reflexivity:} for all \(\beta \subseteq \alpha\) holds \(\alpha \to \beta\)
\item \textbf{Transitivity:} \(\alpha \to \beta \land \beta \to \gamma \implies \alpha \gamma\)
\item \textbf{Connectivity:} \(\alpha \to \beta \implies \alpha \gamma \to \beta \gamma\)
\item \textbf{Union:} \(\alpha \to \beta \land \alpha \to \gamma \implies \alpha \to \beta \gamma\)
\item \textbf{Dedomposition:} \(\alpha \to \beta \gamma \implies \alpha \to \beta \land \alpha \to \gamma\)
\item \textbf{Pseudotransivity:} \(\alpha \to \beta \land \beta \gamma \to \delta \implies \alpha \gamma \to \delta\)
\end{itemize}

\subsection{Attribute Closure}

The attribute closure \(AC(FD, x)\)is the set of all attributes reachable via functional dependency inluding 
the starting point.

\textbf{Example:}


\(R = \{A,B,C,D\}\) with \(FD = \{A \to B, B \to C, AB \to CD\}\)

\begin{itemize}
\item \(AC(FD, A) = \{A, AB, ABC, ABCD\} = ABCD\)
\item \(AC(FD, B) = \{B, BC\} = BC\)
\end{itemize}

\subsection{Canonical Overlapping}

We consider two sets of functional dependencies if and only if their 
attribute closures \(FD^+\) are identical.

For a schema \(R\) with a set of functional dependencies \(FD\) we denote \(FC^c\) as 
\textbf{the Canonical Overlapping} of \(FD\) if the following criteria are satisfied: 

\begin{itemize}
\item \(FD^c \cong FD\), \(FD^{c+} = FD^+\)
\item For all \(\alpha \to \beta \in FD^c\) are minimal in the sense that 
  \begin{itemize}
  \item \(\forall a \in \alpha: FD^c \setminus (\alpha \to \beta) \cup (\alpha \setminus a \to \beta) \ne FD^c\)
  \item \(\forall b \in \beta: FD^c \setminus (\alpha \to \beta) \cup (\alpha \to \beta \setminus b ) \ne FD^c\)
  \end{itemize}
\item All left side attributes appear only one time 
\end{itemize}

\subsubsection{Left-Reduction}

For all \(\alpha \to \beta \in FD\) check if there is a redundant \(a \in \alpha\). \(\beta \subseteq AC(FD, \alpha \setminus a)\).
If that is the case then replace \(\alpha \to \beta\) with \(\alpha \setminus a \to \beta\).

\subsubsection{Right-Reduction}

For all \(\alpha \to \beta \in FD\) check if there is a redundant \(b \in \beta\). 
\(\beta \subseteq AC(FD \setminus \{\alpha \to \beta \} \cup \{\alpha \to \beta \setminus b, \alpha\})\).
If that is the case then replace \(\alpha \to \beta\) with \(\alpha \to \beta \setminus b\).

\subsubsection{Elimination}

Remove \(\alpha \to \emptyset\).

\subsubsection{Union}

If \(\alpha \to \beta_1\) and \(\alpha \to \beta_2\) then combine them into \(\alpha \to \beta_1, \beta_2\).

\textbf{Example:}


Given are \(R=\{ A, B, C, D, E, F \}, FD=\{ ABC \to DEF,A \to BC,D \to EF\} \).

\textbf{Left-Reduction:}

Firt iteration \(ABC \to DEF\)

\begin{itemize}
\item Is \(A\) redundant? No: \(AC(F, BC) = BC\)
\item Is \(B\) redundant? Yes: \(AC(F, AC) = AC \implies ABC \implies ABCDEF\).  
Thus we have \(AC \to DEF, A \to BC, D \to EF\).
\item Is \(C\) redundant? Yes: \(AC(F, A) = A \implies ABC \implies ABCDEF\).  
Thus we have \(A \to DEF, A \to BC, D \to EF\).
\end{itemize}

\textbf{Right-Reduction}

Using \(\{A \to DEF, A \to BC, D \to EF\}\) we proceed with the first dependency 
\(A \to DEF\).

\begin{itemize}
\item Is \(D\) redundant? No, without it we cant reach \(EF\).
\item Is \(E\) redundant? Yes: \(AC(A \to DF, A \to BC, A \to EF, A) = A \implies ADF \implies ADEFBC\).  
Thus we have \(F = \{A \to DF, A \to BC, D \to EF\}\).
\item Is \(F\) redundant? Yes: \(AC(A \to D, A \to BC, D \to EF, A) = A \implies ABC \implies ABCDEF\).  
Thus we have \(F = \{A \to D, A \to BC, D \to EF\}\).
\end{itemize}

\textbf{Elimination and Union}

\[
    FC^c = \{A \to BCD, D \to EF\}
\]

\textbf{Important Notes}

\begin{itemize}
\item By the left-reduction we only care about fd's which have at least two attributes in the right side.
\item By the right-reduction we only care about the fd's with one attribute at the right.
\end{itemize}

\subsection{How to find a Candidate Key}

Given a set of attributes \(\{A,B,C,D\}\) with the following functional dependencies:

\[
A \to B, C \to B, A \to D
\]

The set is already left-reduced. 

\begin{itemize}
\item \(AC(F, A) = ABD\)
\item \(AC(F, C) = CB\)
\end{itemize}

We get \(A \to BD, C \to B\).

For the right reduction: 

\begin{itemize}
\item Is \(B\) redundant? Yes
\end{itemize}

Finally, we get our candidate key \(AC \to ABCD\).

\subsection{Losslessness and Maintenance of Dependency}

Given a relation which needs to be splited into more relations:

\begin{itemize}

    \item \textbf{Losslessness:} After the split, the information in the original relation needs be accesible via joins.

    \item \textbf{Maintenance of Dependency:} The original functional dependencies do not change after the split.

\end{itemize}

\subsection{Normal Forms}

A \textbf{normalized} table are protected from data inconsistencies and make the tables more readable.

\begin{itemize}
\item \textbf{I Normal Form:}
  \begin{itemize}
  \item No row order for information.
  \item No Mixing data types in the same column.
  \item No Tables without a primary key.
  \item No Repeating groups.
  \end{itemize}
\end{itemize}

This normal is mostly given in relational algebra problems.

\begin{itemize}
\item \textbf{II Normal Form:} Each non-key attribute must depend on the entire primary key.
  \begin{itemize}
  \item The relation is in the second normal form if it is already in the first normal form and each \textbf{non-key} attribute is functionally dependent on the \textbf{key-candidate}.
  \end{itemize}
\end{itemize}

\textbf{Example I:}

\(R = \{A,B,C,D\}, \quad FD = \{A \to BC, B \to A\}\) not in 2NF due to \(D\) not being present.

\textbf{Example II:}

\(R = \{A,B,C,D\}, \quad FD = \{A \to B, B \to A, AD \to C\}\)

Candidates \(AD, BD\), not prime attribute \(C\). Thus, it obeys the 2NF.

\begin{itemize}
\item \textbf{III Normal Form:} Every non-key attribute in a table should depend on the key, the whole key, and nothing but the key.
  \begin{itemize}
  \item The relation is in the third normal form if it is already in the second normal form and each \textbf{non-key} attribute is not dependent on other \textbf{non-key} attribute.
  \item \textbf{Boyce-Cood III Normal Form:} Every attribute in a table should depend on the key, the whole key, and nothing but the key.
  \end{itemize}
\end{itemize}

\textbf{Formal Definition 3NF:}

A schema \(R\) is in the 3NF if for all \( \alpha \to \beta \in FD\) at least one of the following statements 
is true: 

\begin{itemize}
\item \(\alpha \to \beta\) is trivial \(\beta \subseteq \alpha\)
\item \(\beta\) has only prime attributes 
\item \(\alpha\) is a superkey 
\end{itemize}

\textbf{Formal Definition BCNF:}

A schema \(R\) is in the BCNF if for all \( \alpha \to \beta \in FD\) at least one of the following statements 
is true: 

\begin{itemize}
\item \(\alpha \to \beta\) is trivial \(\beta \subseteq \alpha\)
\item \(\alpha\) is a superkey 
\end{itemize}

\textbf{Example I}

\(R = \{A,B,C,D\}, \quad FD = \{A \to BC, C \to D\}\), not in 3NF due to \(D\) being dependent on \(C\).

\begin{itemize}
\item \textbf{IV Normal Form:} Multivalued dependencies in a table must be multivalued dependencies on the key.
\item \textbf{V Normal Form:} Our table in the forth normal form cannot be describable as the logical result of joining some other tables.
\end{itemize}

\subsection{Synthesis Algorithm}

The \textbf{synthesis algorithm} is used to convert relations to the 3NF.

\begin{itemize}
    
    \item Find the Canonical Overlapping.
   
    \item For all \(\alpha \to \beta \in FD^c\) we get a schema \(R_k = \alpha \cup \beta\) in which 
    the attributes of functional dependencies which are in \(R_k\) will be ordered to \(R_k\).
  
    \item A key candidate has to be in \(R_k\). Else, we create another relation with the attributes of 
    the key candidate.
 
    \item Individual \(R_k\) can contain other \(R_m\). Those schemata are not necessary and can be removed.

\end{itemize}

\textbf{Example:}

Given \(R = \{A,B,C,D,E,F\}\) with \(FD = \{A \to A, A \to B, C \to DEF, E \to F\}\).

\textbf{Left-Reduction}

Nothing to do.

\textbf{Right Reduction}

\(A \to \emptyset\) (trivial)  
\(A \to B\) nothing to do.  
\(C \to DEF\) discard \(F\) due to \(C \to DE \implies C \to DEF\).  
\(E \to F\) nothing to do.

Result: \(FC^c = \{A \to B, C \to DE, E \to F\}\).

\textbf{Schemata}

\(\mathcal{R}_{AB}\) with \(F_{AB} = \{A \to B\}\).  

\(\mathcal{R}_{CDE}\) with \(F_{CDE} = \{C \to DE\}\).  

\(\mathcal{R}_{EF}\) with \(F_{EF} = \{E \to F\}\).

\textbf{Test for key candidates}

No relation has the key-candidate \(AC\).

\subsection{Decomposition Algorithm for BCNF}

The \textbf{decomposition algorithm} is used to convert relations to the BCNF. 

A schema \(R\) is in the BCNF if for all \( \alpha \to \beta \in FD\) at least one of the following statements
is true:

\begin{itemize}

    \item \(\alpha \to \beta\) is trivial \(\beta \subseteq \alpha\)

    \item \(\alpha\) is a superkey 

\end{itemize}

\subsubsection*{Algorithm}

\begin{itemize}
    
    \item Let \(Z = \{R\}\).
    
    \item As long as there schemata \(R_i \in Z\) gives which are not in BCNF proceed to 
    
    \begin{itemize}
        \item Split \(R_i\) into \(R_{i1} = \alpha \cup \beta\) and \(R_{i2} = R_i - \beta\).
        \item Remove \(R_i\) from \(Z\) and add \(R_{i1}\) and \(R_{i2}\).
    \end{itemize}

    \item The algorithm works recursively on the subschemas.

\end{itemize}

\textbf{Example:}

\(R = \{A,B,C,D\}\) with \(F = \{BC \to D, AB \to A\}\).

\(AB \to A\) is in BCNF due to \(AB\) being a key-candidate and \(A \subseteq AB\)

On the other hand \(BC \to D\) is not in BCNF due to \(BC\) not being a key-candidate. 

Thus, we split \(R\) into

\(F_{R_1} = \{BC \to D\}\) and \(F_{R_2} = \{AB \to A\}\).

\subsection{Optimization}

The following equivalences hold for relational algebra expressions: 

\begin{itemize}
    
    \item \(\sigma C_1 \land C_2 \land \dots \land C_n \cong \sigma_{C_1} ( \sigma_{C_2} ( \dots ( \sigma_{C_n} (R)) \dots ))\)
    
    \item \(\sigma_{C_1} (\sigma_{C_2} (R)) \cong \sigma_{C_1} (\sigma_{C_2} (R))\)
    
    \item If \(L_1 \subseteq L_2 \subseteq \dots \subseteq L_n\) then 
        \[
            \pi_{L_1} (\pi_{L_2} ( \dots ( \pi_{L_n} (R)) \dots ) ) \cong \pi_{L_1} (R) 
        \]
    
    \item \(\pi_{A_1, \dots, A_n} (\sigma_{C} (R)) \cong \sigma_C ( \pi_{A_1, \dots, A_n}(R))\)
    
    \item \(\times, \cup, \cap\) and \(\bowtie\) are commutative
    
    \item Pushing Selections, select as early as possible \( \sigma_C \left(R \bowtie_{C_J} S \right) \cong \sigma_C (R) \bowtie_{C_J} S\)

    Special Case: 

    \[
     \sigma_{C_1 \land C_2} \left(R \bowtie_{C_J} S \right) \cong \sigma_{C_1} (R) \bowtie_{C_J} \sigma_{C_2} (S)
    \]

    \item \(\pi_L (R \bowtie_{C_J} S) \cong \pi_{A_1, \dots, A_n}(R) \bowtie_{C_J} \pi_{b_1, \dots, b_n}(S)\)
    
    \item \(\times, \cup, \cap\) and \(\bowtie\) are associative
   
    \item \(\sigma\) is distributive with \(\cup, \cap, -\) 

    \[
        \sigma_C (R \Theta S) = \sigma_C (R) \Theta \sigma_C (S)
    \]
    
    \item \(\pi\) is distributive with \(\cup\) 

    \[
        \pi_C (R \cup S) = \pi_C (R) \cup \pi_C (S)
    \]
    
    \item De Morgam for Joins with selections 

    \[
        \neg (A \land B) = \neg A \lor \neg B \quad \neg (A \lor B) = \neg A \land \neg B
    \]
    
    \item \(\sigma_C (R \times S) \cong R \bowtie_C S\)

\end{itemize}
